\section{Methodology}
\label{sec-model}
In this section, we present the \textbf{M}ulti-\textbf{D}istribution \textbf{L}earning (\textbf{MDL}) framework, a unified architecture designed for multi-scenario learning and multi-task learning in recommendations.
The core idea of our approach is "Tokenize-and-Interact".
As shown in Figure~\ref{fig:model-overall}, our approach mainly consists of two parts: \emph{Unified Information Tokenization} module and \emph{Domain-aware All-Token Interaction} module.
The first module is to transform heterogeneous inputs into a set of discrete latent tokens, \ie feature tokens, scenario tokens, and task tokens.
The second module describes how these tokens interact for achieving unified and effective modeling across multiple scenarios and tasks.
Next, we describe each part in detail.

\subsection{Unified Information Tokenization}
Firstly, we introduce how to represent feature, task and scenario information in a tokenized form uniformly.

\subsubsection{Feature Tokenization}
The input feature is the fundamental of a reommendation model.
In real-world industrial recommendation systems, input features typically include user features $x_u$ (\eg user id, user profile), item features $x_i$ (\eg video id, author id), sequence features $x_{seq}$ (\eg user history click behaviors), and cross features $x_{cross}$. 
Most user features, item features and cross features are first transformed into high-dimensional embeddings $\bm{e}$ through the embedding layers as follows:
\begin{equation}
\bm{e}_{u}, \bm{e}_{i}, \bm{e}_{cross} = \text{EmbLayer}(x_u,x_i,x_{cross}).
\end{equation}
As for sequence features $x_{seq}$, they are usually processed through the designed sequence modules~\cite{DIN,Longer,SIM} to extract user interest and convert them into embeddings $\bm{e}_{seq}$. 
Ultimately, all features are transformed into feature embeddings $\bm{e}_{input} = [\bm{e}_u; \bm{e}_i; \bm{e}_{cross}; \bm{e}_{seq}]$ with varying dimensions.

To achieve complex feature interaction(\eg Self-Attention~\cite{transformer}, MLPMixer~\cite{MLPMixer}), existing large-scale recommendation methods typically transform features into the fixed-number and fixed-dimensional hidden representations $\bm{T}_{f} \in \mathbb{R}^{N_{f} \times d_{f}}$, referred to as feature tokens~\cite{RankMixer,OneTrans,KGBRank}. 
These feature tokens ensure efficient parallel computation and effective feature interaction in later stages.
Following previous work~\cite{RankMixer}, we adopt a semantic-based tokenization approach for generating feature tokens in \modelname.
Specifically, we manually group features into serval semantically coherent clusters with domain knowledge.
These grouped features are sequentially concatenated as the input embeddings $\bm{e}_{input} = [\bm{e}_1; \bm{e}_2; \cdots; \bm{e}_{N_{f}}]$, where $N_{f}$ is the number of feature tokens.
For each grouped feature embedding $\bm{e}_j$, it captures a set of feature embeddings that represent a similar semantic aspect.
Then a projection layer convert the grouped feature embedding $\bm{e}_j$ into a fixed-dimension hidden representation as follows:
\begin{equation}
    \bm{t}_{j} = \text{Proj}(\bm{e}_{j}),
\end{equation}
where $\bm{t}_{j} \in \mathbb{R}^{d_{f}}$, $d_{f}$ is the dimension of the feature token.
Finally, we obtain the feature tokens $\bm{T}_{f} = [\bm{t}_{1};\bm{t}_2;\cdots;\bm{t}_{N_{f}}]$.

\subsubsection{Scenario and Task Tokenization}
Inspired by large language models (LLMs), to further enhance the interaction between feature, scenario,  and task information, we propose tokenizing scenario information and task information to align them with the feature tokens in the same semantic space.
Compared with previous multi-scenario and multi-task recommendation methods, a unified tokenization framework enables multi-scenario and multi-task signals to interact with feature information from the bottom layer and progressively across layers, rather than affecting solely on partial intermediate results or shallow output results.
Similar to how LLMs are applied to specific downstream tasks, these scenario and task tokens act as prompt tokens to activate the model's capabilities for multi-scenario and multi-task recommendation.

Next, we will introduce the tokenization approach for scenarios and tasks.
We begin by taking scenario tokens as an example.
Similar to feature tokens $\bm{T}_{f}$, scenario tokens $\bm{T}_{t} \in \mathbb{R}^{(N_{s} + 1) \times d_{s}}$ are also a fixed number of hidden representations, which are derived from specific feature embeddings through a transformation process.
Specifically, the input feature embeddings for scenario tokens consist of two parts: (1) some important features' extra embeddings $\bm{\hat{e}_{imp}}$ (\eg user id, video id) where is different from the input embeddings of feature tokens $\bm{e}$, but same raw features $x$, and (2) scenario-related prior feature embeddings $\bm{\hat{e}_{spec}}$ (\eg scenario-specific user behavior sequence).
The two components ensure the richness of the initial input information and the differentiation among the different scenario tokens.
Then, the nolinear feed-forward network is used to convert the input embeddings into task tokens as follows:
\begin{equation}
    \bm{t}_{s} = \text{Relu}(\text{FFN}(\bm{\hat{e}_{imp}} \oplus \bm{\hat{e}_{spec}})).
\end{equation}
For each scenario token $\bm{t}_{s}$, the used FFN parameters are different, \ie per-token feed-forward network (Pertoken FFN)~\cite{RankMixer}.
Finally, we obtain the scenario tokens $\bm{T}_{s} = [\bm{t}_{s,1};\bm{t}_{s,2};\cdots;\bm{t}_{s,(N_{s} + 1)}] \in \mathbb{R}^{(N_{s} + 1) \times d_{s}}$, where $d_{s}$ is the dimension of the scenario token, $N_{s}$ is the number of scenario tokens, which equals to the number of recommendation scenarios (\eg homepage feeds or banner).
In addition to the $N_{s}$ scenario tokens, we also construct a global scenario token $\bm{t}_{s,global} \in T_{s}$ that is shared across all scenarios.
The global scenario token is used to learn common knowledge among scenarios.

Similarly, we can obtain task tokens $\bm{T}_{t} \in \mathbb{R}^{N_{t} \times d_{t}}$ by inputting other important feature embeddings and task-related feature embeddings, and using Pertoken FFN transformation.
$N_{t}$ is the number of task tokens, which is equal to the number of prediction tasks of the recommendation model.

\subsection{Domain-aware ALL-Token Interaction}
After obtaining the initial feature tokens, scenario tokens and task tokens, we introduce how to optimize token information by controlling the interaction and information transmission among different tokens, further achieving effective and unified multi-scenario and multi-task recommendation modeling.
The interactions among these tokens can be categorized into three main types: \emph{feature token self-interaction}, \emph{feature-scenario/task token interaction} and \emph{scenario-task token interaction}.
Next, we will introduce in detail.

\subsubsection{Feature Token Self-Interaction}
The interaction between feature tokens serves as a fundamental module in large-scale industrial recommendation models, which has been widely studied in recent works~\cite{OneTrans, RankMixer, hyformer}.
It aims to enhance the model's expressive ability by stacking efficient feature interaction modules to scale model parameters.
For instance, structures like Self-Attention~\cite{transformer} and MLP-Mixer~\cite{MLPMixer}, along with their variants, have been widely applied in the design of feature interaction modules.
The feature interaction layer can be formulated as follows:
\begin{equation}
    \bm{T}_{f}^{(l+1)} = \text{Sef-Interaction}(\bm{T}_{f}^{(l)}),
\end{equation}
where $l$ is the number of layers, $\bm{T}_{f}^{(l)}, \bm{T}_{f}^{(l+1)} \in \mathbb{R}^{N_{t} \times d_{t}}$ represent the feature tokens before and after interaction, respectively.
Specifically, in this work, we follow the feature interaction module design from the previous approach RankMixer~\cite{RankMixer}.
It mainly consists of two parts: TokenMixing module and Pertoken FFN.
So, the feature token self-interaction in \modelname is as follows:
\begin{equation}
\label{eq:feature-interaction}
    \bm{T}_{f}^{(l+1)} = \text{PertokenFFN}(\text{LN}(\text{TokenMixing}(\bm{T}_{f}^{(l)}) + \bm{T}_{f}^{(l)})).
\end{equation}
It is worth noting that this part can be replaced by other feature interaction methods.

For multi-scenario and multi-task learning in recommender system, the interaction among feature tokens provides the foundational modeling capability.
Next, we introduce how the defined scenario tokens and task tokens can be utilized to better activate these capabilities.

\subsubsection{Feature-Scenario/Task Token Interaction}
The interaction between scenario/task tokens and feature tokens primarily aims to enable different scenario tokens or task tokens to utilize distinct feature information, thereby achieving differentiated modeling in multi-scenario and multi-task settings.
Since task and scenario information are tokenized from the bottom layer, they can fully interact with feature information layer by layer and from the bottom up.
This approach alleviates the challenge of insufficient utilization of complex feature interaction modules, which was previously limited to the certain intermediate or output layers.

Specifically, we define a type of interaction called \emph{Domain-aware Attention} between scenario/task tokens and feature tokens.
Without loss of generality, we take the interaction between task tokens $\bm{T}_{t}$ and feature tokens $\bm{T}_{f}$ as an example to illustrate this module.
The \emph{Domain-aware Attention} can be viewed as a variant of cross-attention.
Given the task token $\bm{T}_{t}$, we use it as query, while treating the feature token $\bm{T}_{f}$ as key and value.
Then they are interacted through cross multi-head attention as follows:
\begin{equation}
\label{eq-daa-task}
    \bm{\hat{T}}_t^{(l+1)} = \text{softmax}(\frac{(\bm{W}_Q\bm{T}_t^{(l)})(\bm{W}_K\bm{T}_{f}^{(l)})}{\sqrt{d}})\bm{W}_V\bm{T}_f^{(l)},
\end{equation}
where $l$ is the number of layers, $\bm{W}_Q, \bm{W}_K, \bm{W}_V$ are the QKV projection matrices in the form of a PerToken FFN.
Similarly, for the scenario token $\bm{T}_{s}$, we use it as the query and perform the same multi-head cross-attention interaction with the feature token $\bm{T}_{f}$:
\begin{equation}
\label{eq-daa-scenario}
    \bm{\hat{T}}_s^{(l+1)} = \text{softmax}(\frac{(\bm{W}_Q\bm{T}_s^{(l)})(\bm{W}_K\bm{T}_{f}^{(l)})}{\sqrt{d}})\bm{W}_V\bm{T}_f^{(l)}.
\end{equation}
Through the cross-attention interaction between tokens, specific task tokens and scenario tokens can adaptively and selectively aggregate the rich information from feature tokens based on their own data characteristics, enabling differentiated modeling for different tasks and scenarios. 
Furthermore, since tokenization is performed from the bottom input layer, this interaction can be stacked across multiple layers with feature token self-interaction layer.
The layer-wise cross-attention interaction ensures effective utilization of information within complex structures.

\subsubsection{Scenario-Task Token Interaction}
\label{subsubsec-interaction3}
The interaction between task tokens and scenario tokens primarily focuses on the fusion of their information, enabling differentiated model predictions for specific task within specific scenario in the context of tokenization.
This interaction is mainly controlled by the \emph{Domain-fused Module}.
The core idea is that for each instance, we aim to aggregate only the corresponding scenario token information $\bm{t}_{s}$ for each task token $\bm{t}_{t}$.
This ensures scenario-specific prediction for different recommendation tasks.

Specifically, the \emph{Domain-fused Module} consists of two steps: (1) instance-level scenario token selection, and (2) hybrid information fusion between task and scenario tokens.
Firstly, for each instance in the dataset, we select the corresponding scenario tokens $\{{\bm{t}_{s,i}, \bm{t}_{s,j}}\}$ from the full scenario token set $\bm{T}_s = \{\bm{t}_{s,1}, \bm{t}_{s,2}, \cdots, \bm{t}_{s,{n_s}}, \bm{t}_{s, global}\}$ based on the scenario information associated with each instance, along with the global scenario token $\bm{t}_{s, global}$.
For example, suppose a model is designed to serve $N_{s}$ scenarios. 
For a given instance, it must belong to one or more of these scenarios (in cases where scenario overlap exists). 
Based on this information, the corresponding scenario tokens are selected. 
After selecting the scenario tokens $\{{\bm{t}_{s,i}, \bm{t}_{s,j}}, \bm{t}_{s,global}\}$, we aggregate their information using a mean pooling approach:
\begin{equation}
\label{eq-df1}
    \bm{t}_{s, avg} = \text{MeanPooling}(\{\bm{t}_{s,i}, \bm{t}_{s,j},\bm{t}_{s,global}\}),
\end{equation}
where $\bm{t}_{s, avg}$ represents the aggregated scenario information corresponding to the instance.
Next, we fuse the aggregated scenario token information $\bm{t}_{s, avg}$ into the task token $\bm{T}_{t}$.
For each Task token $\bm{t}_{t,k} \in \bm{T}_{t}$, we directly use a sum pooling approach to integrate the two pieces of information:
\begin{equation}
\label{eq-df2}
    \bm{T}_{t} = \bm{T}_{t} + \bm{t}_{s,avg}.
\end{equation}
In our practice, the simple pooling approach is not only efficient but also sufficiently effective.
Through this tokenized information fusion, the prediction between scenarios and tasks is decoupled, enabling flexible combinations of any scenario and task modeling.

\subsection{MDL Block}
After introducing the main types of interactions in MDL, we now provide a more detailed description of the interaction process within a complete MDL block by using these interactions.

\paratitle{Forward Propagation of Feature Tokens.}
The information propagation of feature tokens primarily relies on their self-interaction.
For each layer of the MDL Block, the propagation of feature tokens is equal to Equation~\ref{eq:feature-interaction}.

\paratitle{Forward Propagation of Scenario Tokens.}
The information propagation of scenario tokens primarily involves domain-aware attention between scenario tokens and feature tokens (Eq~\ref{eq-daa-scenario}), per-scenario FFN, and residual connection.
It can be formulated as follows:
\begin{equation}
    \bm{\hat{T}}_{s}^{(l + 1)} = \text{DomainAwareAttn}(\bm{T}_{s}^{(l)}, \bm{T}_f^{(l + 1)}) + \bm{T}_{s}^{(l)},
\end{equation}
\begin{equation}
    \bm{T}_{s}^{(l + 1)} = \text{PertokenFFN}(\bm{\hat{T}}_{s}^{(l + 1)}) + \bm{\hat{T}}_{s}^{(l + 1)}.
\end{equation}
The per-scenario FFN after the domain-aware attention provides additional nonlinear transformations for each scenario token, enhancing the expressive capability of each scenario token.
And the residual connection ensures the training stability of the scenario tokens and the continuity of information across layers.

\paratitle{Forward Propagation of Task Tokens.}
The information propagation of task tokens primarily consists of domain-aware attention between task tokens and feature tokens (Eq~\ref{eq-daa-task}), domain-fused module (Eq~\ref{eq-df1}, Eq~\ref{eq-df2}), per-task FFN and residual connection.
It can be formulated as follows:
\begin{equation}
    \bm{\hat{T}}_{t}^{(l + 1)} = \text{DomainAwareAttn}(\bm{T}_{t}^{(l)}, \bm{T}_f^{(l + 1)}) + \bm{T}_{t}^{(l)},
\end{equation}
\begin{equation}
    \bm{\tilde{T}}_t^{(l+1)} = \text{DomainFusedModule}(\bm{\hat{T}}_{t}^{(l + 1)}, \bm{\hat{T}}_{s}^{(l + 1)}),
\end{equation}
\begin{equation}
    \bm{T}_t^{(l+1)} = \text{PertokenFFN}(\bm{\tilde{T}}_t^{(l+1)}) + \bm{\tilde{T}}_t^{(l+1)}.
\end{equation}
The pertoken FFN and residual connection are also used in the information propagation of task tokens.

Finally, by stacking $L$ layers of MDL Blocks, we obtain the final output task tokens $\bm{T}_t^{(L)}$.
We attach the logits layer to the final layer of task tokens to serve as the output for the model's different prediction task as follows:
\begin{equation}
    \hat{y}_{n} = \text{LogitsLayer}(\bm{t}_{t,n}^{(L)}),
\end{equation}
where $n \in N_{t}$ is the index of the prediction task.
Since the final output task tokens have already aggregated the corresponding scenario token information for each instance through the domain-fused module, the gradients of instances from a specific scenario will only update the corresponding scenario tokens. 
In this way, scenario tokens are influenced solely by their own scenario information, allowing them to aggregate information from feature tokens and effectively model the differences between scenarios.
And because the tokenized information aggregation is performed bottom-up, the number of final prediction outputs is independent of the scenarios, yet it enables differentiated predictions across scenarios.
